\section{Appendix A: Prompt Strategies}

\subsection{S1 (Neutrosophic)}

\textbf{System}: You are an expert in Neutrosophic Logic. You evaluate
statements using three INDEPENDENT dimensions: Truth (T), Indeterminacy
(I), and Falsity (F), each on {[}0.0, 1.0{]}. These dimensions are NOT
constrained to sum to 1.0. A statement can be simultaneously partially
true AND partially false AND partially indeterminate. Respond with ONLY
a JSON object, no other text.

\textbf{User}: Evaluate this statement on three independent dimensions:

Statement: ``\{statement\}''

\begin{itemize}
\tightlist
\item
  Truth (T): To what degree is this statement true? {[}0.0 to 1.0{]}
\item
  Indeterminacy (I): To what degree is the truth value unknown,
  undetermined, or inherently uncertain? {[}0.0 to 1.0{]}
\item
  Falsity (F): To what degree is this statement false? {[}0.0 to 1.0{]}
\end{itemize}

T, I, and F are independent. They need NOT sum to 1.0.

Respond with ONLY: \{``T'': , ``I'': , ``F'': \}

\subsection{S2 (Probabilistic)}

\textbf{System}: You are a probabilistic classifier. You assign
probabilities to three mutually exclusive categories that MUST sum to
exactly 1.0. Respond with ONLY a JSON object, no other text.

\textbf{User}: Classify this statement into three mutually exclusive
categories whose probabilities sum to 1.0:

Statement: ``\{statement\}''

\begin{itemize}
\tightlist
\item
  T (True): Probability the statement is true
\item
  I (Uncertain): Probability the truth value is unknown or undetermined
\item
  F (False): Probability the statement is false
\end{itemize}

CONSTRAINT: T + I + F must equal 1.0

Respond with ONLY: \{``T'': , ``I'': , ``F'': \}

\subsection{S3 (Entropy-Derived)}

\textbf{System}: You are a binary truth estimator. You estimate the
probability that a statement is true (YES) versus false (NO). The two
probabilities must sum to 1.0. Respond with ONLY a JSON object, no other
text.

\textbf{User}: Estimate the probability that this statement is true
versus false:

Statement: ``\{statement\}''

\begin{itemize}
\tightlist
\item
  P\_yes: Probability the statement is true {[}0.0 to 1.0{]}
\item
  P\_no: Probability the statement is false {[}0.0 to 1.0{]}
\end{itemize}

CONSTRAINT: P\_yes + P\_no must equal 1.0

Respond with ONLY: \{``P\_yes'': , ``P\_no'': \}

\emph{Indeterminacy derived from Shannon entropy:} I = -(p$\cdot$log$_{2}$(p) +
(1-p)$\cdot$log$_{2}$(1-p)) where p = P\_yes

\subsection{S4 (Tensor - Declared Losses)}

\textbf{System}: You are an expert in Neutrosophic Logic and epistemic
honesty. You evaluate statements using three INDEPENDENT dimensions:
Truth (T), Indeterminacy (I), and Falsity (F), each on {[}0.0, 1.0{]}.
These dimensions are NOT constrained to sum to 1.0. Crucially, you must
also declare your LOSSES: what you cannot evaluate, what limits your
assessment, and why your indeterminacy value is what it is. Respond with
ONLY a JSON object, no other text.

\textbf{User}: Evaluate this statement on three independent dimensions,
and declare what you cannot evaluate:

Statement: ``\{statement\}''

\begin{itemize}
\tightlist
\item
  Truth (T): To what degree is this statement true? {[}0.0 to 1.0{]}
\item
  Indeterminacy (I): To what degree is the truth value unknown,
  undetermined, or inherently uncertain? {[}0.0 to 1.0{]}
\item
  Falsity (F): To what degree is this statement false? {[}0.0 to 1.0{]}
\item
  losses: A list of objects, each with:

  \begin{itemize}
  \tightlist
  \item
    ``what'': What you cannot evaluate (brief description)
  \item
    ``why'': Why this limits your assessment
  \item
    ``severity'': How much this affects your evaluation {[}0.0 to 1.0{]}
  \end{itemize}
\end{itemize}

T, I, and F are independent. They need NOT sum to 1.0. You MUST declare
at least one loss. Honesty about limits is required.

Respond with ONLY: \{``T'': , ``I'': , ``F'': , ``losses'':
{[}\{``what'': ``'', ``why'': ``'', ``severity'': \}, \ldots{]}\}

\section{Appendix B: Test Stimuli}

The five test stimuli are from Leyva-Vázquez and Smarandache (2025):

\begin{enumerate}
\def\labelenumi{\arabic{enumi}.}
\tightlist
\item
  \textbf{Paradox (Logical)}: ``This sentence is false.''
\item
  \textbf{Ignorance (Epistemic)}: ``The number of stars in the universe
  is even.''
\item
  \textbf{Vagueness (Fuzzy)}: ``John is 1.75 meters tall, therefore John
  is tall.''
\item
  \textbf{Contradiction (Ethical)}: ``Lying to save an innocent life is
  morally right and wrong at the same time.''
\item
  \textbf{Contingency (Future)}: ``It will rain in New York tomorrow.''
\end{enumerate}

\section{Appendix C: Data Availability}

All code, prompts, and data are available at:
https://github.com/fsgeek/neutrosophic-llm-logic

\begin{itemize}
\tightlist
\item
  \texttt{src/prompts.py}: All four prompt strategies
\item
  \texttt{src/experiment.py}: Experiment runner with strategy selection
\item
  \texttt{data/cross\_vendor\_results.csv}: S1--S3 production data (375
  evaluations)
\item
  \texttt{data/s4\_tensor\_results.csv}: S4 tensor data (125
  evaluations)
\item
  \texttt{data/s4\_mistral\_rerun.csv}: Mistral S4 rerun at 1500
  max\_tokens (25 evaluations)
\item
  \texttt{data/tautology\_s4\_results.json}: Tautology control data (27
  evaluations, Section~\ref{sec:tautology})
\item
  \texttt{data/s5\_ablation\_results.json}: Prompt ablation data (45
  evaluations, Section~\ref{sec:ablation})
\item
  \texttt{data/openai\_neutrosophic\_results.csv}: Original Smarandache
  data
\end{itemize}
